\documentclass[10pt,a4paper]{amsart}
\usepackage[margin=19mm]{geometry}
\usepackage{amsmath,amssymb,mathtools}
\usepackage[scr=rsfs,cal=euler]{mathalfa}
\usepackage{xfrac}
\usepackage[unicode,hidelinks,bookmarks=false]{hyperref}
\newcommand{\ie}{i.e.\ }
\newcommand{\cf}{cf.\ }
\newcommand{\resp}{respectively\ }
\newcommand{\Subsection}{\subsection}
\providecommand{\nobreakdash}{\nobreak\hbox{-}}
\setlength{\emergencystretch}{2em}
\multlinegap0pt
\usepackage{algorithm}
\usepackage{algpseudocode}
\usepackage{listings}
\usepackage{textcomp}
\usepackage{url}
\usepackage{epsfig}
\usepackage{subfig}
\newtheorem{theorem}{Theorem}[section]
\newtheorem{lemma}[theorem]{Lemma}
\newtheorem{proposition}[theorem]{Proposition}
\newtheorem{corollary}[theorem]{Corollary}
\newtheorem{definition}[theorem]{Definition}
\newtheorem{remark}[theorem]{Remark}
\newtheorem{assumption}[theorem]{Assumption}
\usepackage{mathtools}
\usepackage{float}
\usepackage{tikz}
\usepackage{url}
\usepackage{pgfplots}
\usetikzlibrary{matrix,arrows,calc,positioning,shapes,decorations.pathreplacing}
\usepackage{xy}
\newcommand{\lform}[2]{{\big( {#1} \big|\, {#2}\big)}}
\newcommand{\Lform}[2]{{\Big( {#1} \Big|\, {#2}\Big)}}
\newcommand{\scp}[2]{{\left\langle {#1}\, , \, {#2}\right\rangle}}
\newcommand{\dscp}[2]{{\left\langle \!\left\langle {#1}\, , \, {#2}\right\rangle \!\right\rangle}}
\usepackage{bm}
\newcommand{\vect}[1]{\bm{#1}}
\newcommand{\matr}[1]{\bm{#1}}
\newcommand{\quat}[1]{\overline{\vect{#1}}}
\DeclareMathOperator{\conj}{J_\mathbb{H}}
\renewcommand{\S}{\matr{S}} % original: \S is section sign (§)
\renewcommand{\S}{\matr{S}} % original: \S is section sign (§)
\newcommand{\I}{\matr{I}}
\newcommand{\K}{\matr{K}}
\newcommand{\DQ}[1]{\vect{#1}}
\newcommand{\Q}{\DQ{Q}}
\newcommand{\dQ}{\DQ{\delta Q}}
\newcommand{\PrincipalPart}{\mathcal{P}}
\newcommand{\DualPart}{\mathcal{D}}
\newcommand{\Conj}{\mathcal{J}}
\newcommand{\qprod}[2]{#1~\circ~#2}
\newcommand{\npdot}[4][\qprod]{#1{#1{#2}{#3}}{#4}}
\newcommand{\dqprod}[2]{\qprod{#1}{#2}}
\newcommand{\qrot}[2]{\qprod{\qprod{#2^*}{#1}}{#2}}
\newcommand{\qnorm}[1]{\norm{#1}}
\DeclareMathOperator{\sign}{sgn}
\DeclareMathOperator{\atantwo}{atan2}
\newcommand{\R}{\mathbb{R}}
\newcommand{\Rn}{\mathbb{R}^n}
\newcommand{\N}{\mathbb{N}}
\newcommand{\C}{\mathcal{C}}
\newcommand{\D}{\mathcal{D}}
\newcommand{\X}{\mathcal{X}}
\newcommand{\U}{\mathcal{U}}
\newcommand{\Verts}[1]{{\left\Vert #1 \right\Vert}}
\DeclareMathOperator{\diag}{diag}
\DeclareMathOperator{\var}{var}
\DeclareMathOperator{\mean}{\mu}
\DeclareMathOperator{\Var}{\Sigma}
\DeclareMathOperator{\Mean}{\bm\mu}
\newcommand{\ev}[1]{\operatorname{E}\left[{#1}\right]}
\DeclareMathOperator{\rank}{rank}
\DeclareMathOperator{\eig}{eig}
\newcommand{\argmax}{\arg\!\max}
\DeclareMathOperator*{\argmin}{argmin}
\DeclareMathOperator{\tr}{Tr}
\DeclareMathOperator{\Prob}{P}
\newtheorem{assum}{Assumption}
\newtheorem{prop}{Proposition}
\newtheorem{cor}{Corollary}
\newtheorem{rem}{Remark}
\newtheorem{exam}{Example}
\newtheorem{problem}{Problem}
\newcommand{\GP}{\mathcal{GP}}
\newcommand{\z}{\bm z}
\newcommand{\bma}{\bm a}
\newcommand{\x}{\bm x}
\newcommand{\tx}{\tilde{\bm x}}
\newcommand{\m}{\bm m}
\newcommand{\dx}{\dot{\bm x}}
\newcommand{\q}{\bm q}
\newcommand{\dq}{\dot{\bm q}}
\newcommand{\e}{\bm e}
\newcommand{\de}{\dot{\bm e}}
\newcommand{\ddq}{\ddot{\bm q}}
\newcommand{\f}{\bm{f}}
\newcommand{\g}{\bm{g}}
\renewcommand{\u}{\bm{u}}
\newcommand{\y}{\bm{y}}
\newcommand{\leig}{\underline{\lambda}}
\newcommand{\geig}{\bar{\lambda}}
\newcommand{\ubar}[1]{\underline{#1}}
\usepackage{mathtools}
\newcommand{\M}{\matr{M}}
\newcommand{\Ve}{U_{ext}^{ij}}
\newcommand{\Vezero}{U_{i,0}^{\text{ext}}}
\newcommand{\Ad}{\text{Ad}}
\newcommand{\ad}{\text{ad}}
\newcommand{\se}{\mathfrak{se}(2)}
\DeclareMathOperator{\spn}{span}
\begin{document}
\title[Feasibility and Singularity in High-Order Safety-Critical Co]{Feasibility and Singularity in High-Order Safety-Critical Control for Quadrotor UAVs}
\author{Omayra Yago Nieto, Leonardo Colombo}
\date{}
\maketitle
\noindent\emph{Source and licence.} Feasibility and Singularity in High-Order Safety-Critical Control for Quadrotor UAVs. Version arXiv:2609.19362v1 (CC BY 4.0). This material is distributed under the Creative Commons Attribution 4.0 International licence, http://creativecommons.org/licenses/by/4.0/, as recorded in the arXiv OAI licence metadata for this exact version and shown on the abstract page https://arxiv.org/abs/2609.19362v1. Full rights notice: licenses/arxiv-2609.19362-CC-BY-4.0.txt.\par\smallskip
\noindent\textit{Editorial scope.} Counted unit: the original \texttt{theorem} environment \texttt{thm:singularity\_removal} at source lines 674--707 together with its complete proof at lines 709--772; the delivered document keeps source lines 323--773 so that every referenced label and every citation resolves inside it. Native editable source is the paper's own concatenated TeX; the AMS wrapper is editorial formatting only. No publisher class, upstream script or unrelated artwork is redistributed.
\par\smallskip\noindent\textit{Reference note.} This article ships no compiled bibliography (biblatex); the entries delivered here are composed from the author's own BibTeX (.bib) fields, with no field invented and no entry dropped.
\section{Thrust-Based High-Order Safety and Aggregate Feasibility}\label{secII}

\subsection{Pairwise HOCBF safety and thrust effectiveness}

Consider $N$ quadrotor UAVs indexed by
$\mathcal V=\{1,\ldots,N\}$. The state of agent $i$ is
$x_i=(p_i,v_i,R_i,\Omega_i)$, with $p_i,v_i,\Omega_i\in
\mathbb R^3$ and $R_i\in SO(3)$, and its control input is
$(f_i,\tau_i)\in\mathbb R\times\mathbb R^3$. The dynamics are (see \cite{lee2010geometric}, for instance)
\begin{align}
    \dot p_i &=v_i,\qquad
    \dot R_i=R_i\widehat{\Omega}_i,                         \label{eq:quad_kin}\\
    m_i\dot v_i
    &=f_iR_ie_3+m_i g+\Delta_{f,i}(p_i,v_i),                    \label{eq:quad_trans}\\
    J_i\dot\Omega_i+\Omega_i\times J_i\Omega_i
    &=\tau_i+\Delta_{\tau,i}(x_i),                          \label{eq:quad_rot}
\end{align}
where $m_i>0$, $J_i\succ0$, $e_3=(0,0,1)^\top$,
$g=-g_0e_3$, \(g_0>0\) and $\widehat a\,b=a\times b$.

Throughout, the analysis is restricted to a compact operating
set. The force perturbations $\Delta_{f,i}(p_i,v_i)$ are
$C^2$, while $\Delta_{\tau,i}(x_i)$ are locally Lipschitz
there. The restricted dependence
$\Delta_{f,i}=\Delta_{f,i}(p_i,v_i)$ ensures that, after the
dynamic extension introduced in Section~\ref{sec:singularity_free}, the uncertain
fourth-order contribution modifies only the drift term and not
the extended-input coefficient. We consider locally Lipschitz feedback laws. Under
these assumptions, the closed-loop solutions are unique and
the derivatives entering the fourth-order HOCBF are well
defined classically along trajectories. The fourth-order construction is required because the attitude
torques first appear explicitly after four differentiations of
the pairwise distance barrier.

For the safety analysis, let $\mu_{f,i}(p_i,v_i)$ be a learned
estimate of $\Delta_{f,i}(p_i,v_i)$ obtained here from GP
regression~\cite{rasmussen2006gpml}, and define
$\bar\Delta_i:=\Delta_{f,i}-\mu_{f,i}$. Let $\bar d_{f,i}(p_i,v_i)$ denote the associated
confidence bound, so that $\|d_{f,i}(p_i,v_i)\|\le
\bar\Delta_i(p_i,v_i)$ on the operating set.

Collision avoidance is encoded by an undirected graph
$\mathcal G=(\mathcal V,\mathcal E)$. For $(i,j)\in\mathcal E$,
define
\begin{equation}
    h_{ij}:=\|p_i-p_j\|^2-d_{\min}^2,
    \qquad
    \mathcal C_{ij}:=\{x:h_{ij}(x)\ge0\},
    \label{eq:hij}
\end{equation}
where $x=\operatorname{col}(x_1,\ldots,x_N)$ and
$d_{\min}>0$.

Because the singularity studied below corresponds to a loss of uniform relative degree, we briefly recall the
least-relative-degree HOCBF construction of~\cite{tan2022highorder}. For
$\dot x=F(x)+G(x)u$, a function $h\in C^r$ has least
relative degree $r$ on $D$ if
$L_GL_F^kh=0$, $k=0,\ldots,r-2$, where $L_F$ denotes
the Lie derivative along $F$ and
$L_Gq:=\nabla q^\top G$. It has uniform relative degree
$r$ if also $L_GL_F^{r-1}h\neq0$ on $D$. With sufficiently differentiable extended class-$\mathcal K$
functions $\alpha_k$,
$\varphi_0:=h$ and
$\varphi_k:=\dot\varphi_{k-1}+\alpha_k(\varphi_{k-1})$,
a locally Lipschitz feedback satisfying $\varphi_r\ge0$
renders
$\bigcap_{k=0}^{r-1}\{\varphi_k\ge0\}$ forward invariant
whenever the HOCBF condition is feasible. 


 Let $p_{ij}:=p_i-p_j$, $v_{ij}:=v_i-v_j$, and $f_{ij}:=\operatorname{col}(f_i,f_j)$. Since $\dot h_{ij}=2p_{ij}^\top v_{ij}$, the true second-order HOCBF defined by $\varphi_{0,ij}:=h_{ij}$, $\varphi_{1,ij}:=\dot h_{ij}+\alpha_1(h_{ij})$, and $\varphi_{2,ij}:=\dot\varphi_{1,ij}+\alpha_2(\varphi_{1,ij})$ admits the decomposition \begin{equation} \varphi_{2,ij} = \widehat c_{ij}(x) +B_{ij}(x)f_{ij} +\delta_{ij}^{(2)}(x), \label{eq:pairwise_hocbf_uncertain} \end{equation} where \begin{equation} B_{ij}(x) := 2\begin{bmatrix} \dfrac{p_{ij}^\top R_ie_3}{m_i} & -\dfrac{p_{ij}^\top R_je_3}{m_j} \end{bmatrix}, \label{eq:Bij} \end{equation} and $\widehat c_{ij} :=2\|v_{ij}\|^2 +2p_{ij}^\top(\mu_{f,i}/m_i-\mu_{f,j}/m_j) +\alpha_1'(h_{ij})\dot h_{ij} +\alpha_2(\dot h_{ij}+\alpha_1(h_{ij}))$. Moreover, $\delta_{ij}^{(2)} :=2p_{ij}^\top(d_{f,i}/m_i-d_{f,j}/m_j)$ satisfies \[ \left|\delta_{ij}^{(2)}\right|\leq 2\|p_{ij}\|
\left( \frac{\bar\Delta_i}{m_i} + \frac{\bar\Delta_j}{m_j}\right)=:\rho_{ij}^{(2)}.\] Here, $\widehat c_{ij}$ contains the known terms and the
GP-based force estimates, whereas $\delta_{ij}^{(2)}$ collects
the residual force uncertainty. Writing the input-dependent
term separately as $B_{ij}f_{ij}$ isolates the relative-thrust
coefficient whose norm is used below to quantify thrust
effectiveness. Hence a robust pairwise HOCBF condition is $B_{ij}(x)f_{ij}\ge-\widehat c_{ij}(x)+\rho_{ij}^{(2)}(x)$. 
 
 \begin{definition} The pairwise thrust-effectiveness measure and thrust singular set are $\gamma_{ij}^{\rm th}(x):=\|B_{ij}(x)\|$ and $\Sigma_{ij}^{\rm th}:=\{x:B_{ij}(x)=0\}$. \end{definition} 
 
 Equivalently,
$\Sigma_{ij}^{\rm th}
=\{x:p_{ij}^\top R_ie_3=0,\,
p_{ij}^\top R_je_3=0\}$.
Thus $h_{ij}$ has uniform relative degree two away from
$\Sigma_{ij}^{\rm th}$. At the singular set, the thrust
magnitudes cannot modify $\varphi_{2,ij}$ instantaneously.
The uncertainty margin shifts the feasibility level but leaves
the thrust singular set unchanged.
 
\subsection{Robust aggregate HOCBF feasibility}

Let $f:=\operatorname{col}(f_1,\ldots,f_N)$ and assume bounded
thrust inputs
$\displaystyle{\mathcal F:=\prod_{i=1}^N[f_i^{\min},f_i^{\max}]},\,
 0<f_i^{\min}<f_i^{\max}$.
Embedding $B_{ij}(x)$ into the full thrust vector, let
$\bar B_{ij}(x)\in\mathbb R^{1\times N}$ denote the row whose
only nonzero entries are the $i$th and $j$th components
given by~\eqref{eq:Bij}, and define
$r_{ij}^{\rm rob}(x):=
-\widehat c_{ij}(x)+\rho_{ij}^{(2)}(x)$. The robust pairwise HOCBF condition is therefore
\begin{equation}
    \bar B_{ij}(x)f\ge r_{ij}^{\rm rob}(x),
    \qquad (i,j)\in\mathcal E.
    \label{eq:pairwise_affine_constraint}
\end{equation}

Stacking all pairwise inequalities gives
\begin{equation}
    B(x)f\ge r^{\rm rob}(x),
    \qquad f\in\mathcal F,
    \label{eq:aggregate_constraints}
\end{equation}
where $B(x)\in\mathbb R^{M\times N}$,
$M:=|\mathcal E|$, and $r^{\rm rob}(x)\in\mathbb R^M$
stacks the corresponding $r_{ij}^{\rm rob}$. Robust aggregate
feasibility refers to the simultaneous satisfaction of all
pairwise inequalities by a single $f\in\mathcal F$, i.e.,
their pointwise compatibility under the shared input
constraints~\cite{tan2022compatibility}.

\begin{definition}
For a fixed state $x$, define
\begin{equation}
    \Gamma_{\rm th}^{\rm rob}(x)
    :=
    \max_{f\in\mathcal F}
    \min_{(i,j)\in\mathcal E}
    \bigl(
        \bar B_{ij}(x)f-r_{ij}^{\rm rob}(x)
    \bigr).
    \label{eq:Gamma_th_primal}
\end{equation}
Thus, $\Gamma_{\rm th}^{\rm rob}(x)$ is the largest common
margin achievable across all robust pairwise HOCBF
inequalities with a single admissible thrust vector.
\end{definition}

\begin{remark}
Pairwise thrust effectiveness and robust aggregate feasibility
are distinct properties. The quantity $\gamma_{ij}^{\rm th}$
measures the input effectiveness of a single edge, whereas
$\Gamma_{\rm th}^{\rm rob}$ measures simultaneous feasibility
under the shared bounded thrust inputs and uncertainty margins.
\end{remark}

For $y\in\mathbb R^N$, let
$\sigma_{\mathcal F}(y):=\max_{f\in\mathcal F}y^\top f$
denote the support function of $\mathcal F$~\cite{boyd2004convex}.

\begin{theorem}\label{thm:robust_aggregate_feasibility}
Consider the robust bounded-thrust HOCBF inequalities
\eqref{eq:aggregate_constraints}. Let
$\Delta_M:=\{\lambda\in\mathbb R_{\ge0}^M:
\mathbf 1^\top\lambda=1\}$. Then
\begin{equation}
    \Gamma_{\rm th}^{\rm rob}(x)
    =
    \min_{\lambda\in\Delta_M}
    \left[
        \sigma_{\mathcal F}\big(B(x)^\top\lambda\big)
        -
        \lambda^\top r^{\rm rob}(x)
    \right].
    \label{eq:Gamma_th_dual}
\end{equation}
Consequently, \eqref{eq:aggregate_constraints} is feasible
if and only if $\Gamma_{\rm th}^{\rm rob}(x)\ge0$.
Moreover, $\Gamma_{\rm th}^{\rm rob}(x)>0$ if and only if
there exist $f\in\mathcal F$ and $\varepsilon>0$ such that
$B(x)f\ge r^{\rm rob}(x)+\varepsilon\mathbf 1$.
\end{theorem}

\begin{proof}
For any $s\in\mathbb R^M$, every $\lambda\in\Delta_M$
defines a convex combination of its components, so
$\lambda^\top s\ge\min_k s_k$, with equality at
$\lambda=e_{k^\star}$ for
$k^\star\in\arg\min_k s_k$. Hence
$\min_k s_k
    =
    \min_{\lambda\in\Delta_M}\lambda^\top s$.

Applying this identity to
$s=B(x)f-r^{\rm rob}(x)$ in
\eqref{eq:Gamma_th_primal} gives
\[
    \Gamma_{\rm th}^{\rm rob}(x)
    =
    \max_{f\in\mathcal F}
    \min_{\lambda\in\Delta_M}
    \lambda^\top\big(B(x)f-r^{\rm rob}(x)\big).
\]
The sets $\mathcal F$ and $\Delta_M$ are compact and convex,
and the objective is continuous and affine in each variable.
Hence, by Sion's minimax theorem~\cite{sion1958general},
\[
    \Gamma_{\rm th}^{\rm rob}(x)
    =
    \min_{\lambda\in\Delta_M}
    \max_{f\in\mathcal F}
    \lambda^\top\big(B(x)f-r^{\rm rob}(x)\big).
\]
For fixed $\lambda$,
\[
    \max_{f\in\mathcal F}\lambda^\top B(x)f
    =
    \sigma_{\mathcal F}\big(B(x)^\top\lambda\big),
\]
which proves~\eqref{eq:Gamma_th_dual}.

By~\eqref{eq:Gamma_th_primal},
$\Gamma_{\rm th}^{\rm rob}(x)\ge0$ if and only if there exists
$f\in\mathcal F$ such that
$B(x)f\ge r^{\rm rob}(x)$.
Likewise, if $\Gamma_{\rm th}^{\rm rob}(x)>0$, an optimizer
$f^\star$ satisfies
$B(x)f^\star\ge r^{\rm rob}(x)
+\Gamma_{\rm th}^{\rm rob}(x)\mathbf 1$.
The converse follows directly from
\eqref{eq:Gamma_th_primal}.
\end{proof}

\medskip

\begin{corollary}
Let $\bar f:=\operatorname{col}(\bar f_1,\ldots,\bar f_N)$ and
$d^f:=\operatorname{col}(d_1^f,\ldots,d_N^f)$, where
$\bar f_i:=(f_i^{\max}+f_i^{\min})/2$ and
$d_i^f:=(f_i^{\max}-f_i^{\min})/2$. Then
\begin{equation}
    \Gamma_{\rm th}^{\rm rob}(x)
    =
    \min_{\lambda\in\Delta_M}
    \left[
        \bar f^\top B(x)^\top\lambda
        +
        (d^f)^\top|B(x)^\top\lambda|
        -
        \lambda^\top r^{\rm rob}(x)
    \right].
    \label{eq:Gamma_th_box}
\end{equation}
\end{corollary}

\begin{proof}
Let $y:=B(x)^\top\lambda$. Since $\mathcal F$ is a Cartesian
product of intervals,
\[
    \sigma_{\mathcal F}(y)
    =
    \sum_{i=1}^N
    \max_{f_i\in[f_i^{\min},f_i^{\max}]}y_i f_i
    =
    \bar f^\top y+(d^f)^\top|y|.
\]
Substituting $y=B(x)^\top\lambda$ into
\eqref{eq:Gamma_th_dual} yields~\eqref{eq:Gamma_th_box}.
\end{proof}




\section{Torque-Aware HOCBF Control}
\label{sec:singularity_free}


\subsection{Torque-aware dynamic extension}\label{secIIIB}

To identify the extended input channel without differentiating
the uncertain terms in
\eqref{eq:quad_trans}--\eqref{eq:quad_rot},
consider the reference dynamics obtained by omitting these
terms. Uncertainty is reintroduced at the HOCBF level below.
Let $h_{ij,\mathrm{ref}}^{(k)}$ denote the $k$th time
derivative of $h_{ij}$ along the reference dynamics. This
notation distinguishes reference-dynamics quantities from the
GP-based quantities denoted by hats in Section~\ref{secII}.

Let $\zeta_i:=R_ie_3$. From~\eqref{eq:quad_kin},
$\dot\zeta_i=-R_i\widehat e_3\Omega_i$.
Introduce $q_i:=\dot f_i$ and
$\nu_i:=\dot q_i=\ddot f_i$, and define
$z_i:=(p_i,v_i,R_i,\Omega_i,f_i,q_i)$, $z:=\operatorname{col}(z_1,\ldots,z_N)$, $w_i:=\operatorname{col}(\nu_i,\tau_i)$, and
$w_{ij}:=\operatorname{col}(\nu_i,\nu_j,\tau_i,\tau_j)$.

Since gravity cancels from the relative dynamics, let
$a_{ij}:=f_i\zeta_i/m_i-f_j\zeta_j/m_j$. Then
\[
h_{ij,\mathrm{ref}}^{(3)}
=6v_{ij}^\top a_{ij}+2p_{ij}^\top\dot a_{ij},
\,
h_{ij,\mathrm{ref}}^{(4)}
=6a_{ij}^\top a_{ij}
+8v_{ij}^\top\dot a_{ij}
+2p_{ij}^\top\ddot a_{ij}.
\]
For $a_i:=f_i\zeta_i/m_i$,
\[
\dot a_i
=\frac{q_i}{m_i}\zeta_i+\frac{f_i}{m_i}\dot\zeta_i,
\qquad
\ddot a_i
=\frac{\nu_i}{m_i}\zeta_i
+\frac{2q_i}{m_i}\dot\zeta_i
+\frac{f_i}{m_i}\ddot\zeta_i.
\]
Using the reference rotational dynamics,
$\ddot\zeta_i=d_{\zeta,i}
-R_i\widehat e_3J_i^{-1}\tau_i$, where
$d_{\zeta,i}:=
-R_i\widehat\Omega_i\widehat e_3\Omega_i
+R_i\widehat e_3J_i^{-1}
(\Omega_i\times J_i\Omega_i)$. Consequently,
\begin{equation}
    h_{ij,\mathrm{ref}}^{(4)}
    =
    \ell_{ij,\mathrm{ref}}^{(4)}(z)
    +
    A_{ij}^{\mathrm{ext}}(z)w_{ij}.
    \label{eq:h4_affine}
\end{equation}
where $\ell_{ij,\mathrm{ref}}^{(4)}$ collects the reference drift terms and
\begin{equation}
\begin{aligned}
A_{ij}^{\rm ext}(z)
=
2\Bigg[
&\frac{p_{ij}^{\top}R_ie_3}{m_i},
-\frac{p_{ij}^{\top}R_je_3}{m_j},\\
&-\frac{f_i}{m_i}
 p_{ij}^{\top}R_i\widehat e_3J_i^{-1},
\frac{f_j}{m_j}
 p_{ij}^{\top}R_j\widehat e_3J_j^{-1}
\Bigg].
\end{aligned}
\label{eq:Aext}
\end{equation}
For quantitative input-effectiveness measures, let
$S_{ij}:=\operatorname{diag}
(s_{\nu,i},s_{\nu,j},s_{\tau,i}I_3,s_{\tau,j}I_3)\succ0$
contain fixed characteristic input scales and define
$\gamma_{ij}^{\rm ext,S}(z):=
\|A_{ij}^{\rm ext}(z)S_{ij}\|$.

 We use the same notation $\mathcal C_{ij}$ for its natural lift to the extended state space. We next show that this
extended input channel does not vanish at a thrust singularity.




\begin{theorem}
\label{thm:singularity_removal}
Let $(i,j)\in\mathcal E$ and suppose
$z\in\mathcal C_{ij}$. If at least one of $f_i$ or $f_j$
is nonzero, then $A_{ij}^{\rm ext}(z)\neq0$. In particular,
if $x\in\Sigma_{ij}^{\rm th}$ and $f_i\neq0$, then
\begin{equation}
    \gamma_{ij}^{\rm ext,S}(z)
\ge
\frac{2|f_i|s_{\tau,i}}
{m_i\lambda_{\max}(J_i)}
\|p_{ij}\|>0.
    \label{eq:gamma_ext_bound}
\end{equation}
and analogously for agent $j$. Moreover, on the admissible
thrust domain $f_i\in[f_i^{\min},f_i^{\max}]$,
\begin{equation}
\gamma_{ij}^{\rm ext,S}(z)
\ge
\frac{2d_{\min}}{m_i}
\min\left\{
s_{\nu,i},
\frac{f_i^{\min}s_{\tau,i}}
{\lambda_{\max}(J_i)}
\right\}
>0.
\label{eq:gamma_ext_uniform_bound}
\end{equation}

Consequently,
$h_{ij}$ has uniform relative degree four
with respect to $w_{ij}$ along the reference extended dynamics
whenever at least one vehicle has nonzero thrust.
\end{theorem}

\begin{proof}
If $B_{ij}(x)\neq0$, the first two components of
$A_{ij}^{\rm ext}$ in~\eqref{eq:Aext} cannot both vanish.

Consider $x\in\Sigma_{ij}^{\rm th}$. Since $z\in\mathcal C_{ij}$,
$\|p_{ij}\|\ge d_{\min}>0$.
Let $q_i^p:=R_i^\top p_{ij}$. Then
$\|q_i^p\|=\|p_{ij}\|$ and, by
$x\in\Sigma_{ij}^{\rm th}$,
$(q_i^p)^\top e_3=0$. Hence
$\|(q_i^p)^\top\widehat e_3\|
    =
    \|q_i^p\|
    =
    \|p_{ij}\|$. Since $J_i\succ0$,
\[
\begin{aligned}
\left\|
p_{ij}^\top R_i\widehat e_3J_i^{-1}
\right\|
\ge
\sigma_{\min}(J_i^{-1})
\left\|
(q_i^p)^\top\widehat e_3
\right\| =
\frac{\|p_{ij}\|}{\lambda_{\max}(J_i)}.
\end{aligned}
\]
Thus the $i$th scaled torque block of
$A_{ij}^{\rm ext}S_{ij}$ has norm at least
$\frac{2|f_i|s_{\tau,i}}{m_i\lambda_{\max}(J_i)}
\|p_{ij}\|$, and~\eqref{eq:gamma_ext_bound} follows.
The argument for agent $j$ is identical. 

For the uniform bound, write
$q_i^p=(q_{i,1}^p,q_{i,2}^p,q_{i,3}^p)^\top$.
The $i$th thrust-acceleration and torque blocks give
\[
\begin{aligned}
\|A_{ij}^{\rm ext}S_{ij}\|^2
&\ge
\frac{4}{m_i^2}
\left[
s_{\nu,i}^2(q_{i,3}^p)^2
+
\frac{f_i^2s_{\tau,i}^2}
{\lambda_{\max}^2(J_i)}
\left((q_{i,1}^p)^2+(q_{i,2}^p)^2\right)
\right] \\
&\ge
\frac{4d_{\min}^2}{m_i^2}
\min\left\{
s_{\nu,i}^2,
\frac{(f_i^{\min})^2s_{\tau,i}^2}
{\lambda_{\max}^2(J_i)}
\right\},
\end{aligned}
\]
which proves~\eqref{eq:gamma_ext_uniform_bound}. Since $w_{ij}$
does not enter $h_{ij}^{(k)}$
for $k\le 3$, while the coefficient of $w_{ij}$
in $h_{ij}^{(4)}$
is $A_{ij}^{\rm ext}\neq 0$, the relative-degree claim follows.
\end{proof}

\begin{thebibliography}{99}
\bibitem[]{boyd2004convex} Stephen Boyd; Lieven Vandenberghe. Convex Optimization. (2004)
\bibitem[]{lee2010geometric} Lee, Taeyoung; Leok, Melvin; McClamroch, N Harris. Geometric tracking control of a quadrotor UAV on SE (3). 49th IEEE conference on decision and control (CDC) (2010) 5420--5425
\bibitem[]{rasmussen2006gpml} Carl Edward Rasmussen; Christopher K. I. Williams. Gaussian Processes for Machine Learning. (2006)
\bibitem[]{sion1958general} Maurice Sion. On General Minimax Theorems. Pacific Journal of Mathematics 8 (1958) 171--176
\bibitem[]{tan2022compatibility} Xiao Tan; Dimos V. Dimarogonas. Compatibility Checking of Multiple Control Barrier Functions for Input Constrained Systems. 2022 IEEE 61st Conference on Decision and Control (CDC) (2022) 939--944
\bibitem[]{tan2022highorder} Xiao Tan; Wenceslao Shaw Cortez; Dimos V. Dimarogonas. High-Order Barrier Functions: Robustness, Safety, and Performance-Critical Control. IEEE Transactions on Automatic Control 67 (2022) 3021--3028
\end{thebibliography}
\end{document}
